\subsection{Economic impacts}

% Teixeira2020 té una secció 3.3 Social equity i 3.4 economy 
% També hauria de mirar els apartats 4. Modal shift and car reduction i 5. Synergies with other modes (5.1. PT, 5.2. Cycling promotion)

% literally copied from Qiu2018 pg 3
% Therefore, a bike-sharing service can increase residents’ leisure time. According to economists, leisure time has a real impact on the economy~\cite{Bonilla-Alicea2019}.


% Copiat literal de Bonilla-Alicea2019 pg 139
% BSS can also promote economic growth by increasing real estate devel-opment, helping companies secure more talented workers, and increasing retail visibility and sales (Boland, Murphy, Armstrong, & Barry 2012;Bullock, Brereton, & Bailey, 2017; Urban Land Institute, 2016). BSS can also help companies secure more talented workers by providing a convenient and efficient way to commute to work.

% Literalment de Luo2019 pg 181
% Moreover, manufacturing the bikes, electronic components, and sharing infrastructure requires resource and energy inputs and generates wastes and emissions (Berkhout and Hertin, 2001; Chester and Horvath, 2009; Coelho and Almeida, 2015). Producing too many bikes that exceed the demand could also lead to resource wastage and waste management problems. In Shanghai, more than 1.5 million dock-less bikes were in the street and the retired or abandoned bikes clogged the sidewalks and created huge piles of bike wastes (Benjamin, 2017). 

% The economic dimension of bike-sharing systems encompasses several facets: direct financial performance of the schemes, cost savings for users, external economic benefits (e.g. health care savings, productivity gains), and indirect effects on local economies. While BSS are not typically profitable as standalone enterprises (many require public subsidies or sponsor support), research indicates that their broader economic impact can be positive when benefits are monetized.

% One major economic impact comes via the health benefits described above. Lives saved and illnesses prevented translate into economic value – reduced medical costs, increased productivity from a healthier population, and lower absenteeism. For example, the European analysis by Otero et al. estimated an annual economic benefit of €18 million from the health gains of the 12 cities’ bike-share programs (Otero2018). In the U.S., the \$36 million per year valuation of health benefits by Clockston \& Rojas-Rueda (2021) similarly represents a significant return in public health terms. These figures can be compared to the operating costs of the systems to assess societal return on investment. Murphy and Usher (2015), analyzing Dublin’s bike-share, reported a strongly positive benefit–cost ratio (well above 1) when accounting for health benefits, even under conservative scenarios. In other words, the value of improved health and other external benefits outweighed the city’s expenditures on the program (Murphy2015).

% Beyond health, bike-sharing may confer personal economic savings for users. Those who shift from driving or ride-hailing to bike-share save on fuel, car maintenance, parking or fare costs. Survey data indicates that some low-income users are motivated by the financial savings of cycling versus transit fares or car expenses (Fishman2016). However, because many bike-share trips substitute for walking or are taken by people who already own transit passes, the average direct user saving per trip is often small. More significant are the time savings in certain scenarios: as noted earlier, using bike-share for last-mile connections can shorten commute times (Qiu2018), effectively giving people back free time that has real economic value (often quantified as value of leisure or increased productivity). Qiu and He (2018) calculated that the time saved by Beijing bike-share users in their commutes could increase the city’s GDP by on the order of 1–2 billion yuan annually, by enabling more productive use of time (Qiu2018). While this estimate relies on assumptions and is specific to Beijing’s context, it illustrates the concept that faster trips and improved accessibility have economic merit.

% On a city-wide scale, bike sharing’s presence can also have indirect economic effects. It can boost local retail and tourism by making it easier for people to move around shopping districts or visit attractions. Tourists often use bike-share to explore cities cheaply; for instance, casual usage spikes have been observed near tourist sites in cities like Paris and New York. This can lead to increased spending at local businesses. Moreover, bike-share data (anonymized trip records) have economic value too – cities have used the data to understand consumer movement patterns, aiding urban planning and potentially benefiting commercial location decisions. Another oft-cited impact is the creation of jobs related to bike-share operations: systems require teams for maintenance, bicycle redistribution, customer service, and administration. Although these programs are not massive employers, they do create new jobs in the green transport sector and associated industries (manufacturing bicycles, software development for apps, etc.).

% From a public finance perspective, justifying expenditure on BSS often hinges on the external benefits versus costs. Cost–benefit analyses in several cities have found that when health benefits (and sometimes congestion or environmental benefits) are assigned monetary values, the net benefit is positive (Fishman2016; Murphy2015). This has been a compelling argument for continued public investment in bike-share as a part of sustainable transport strategy. For example, a cost–benefit study of Barcelona’s BSS found that for every euro spent on the program, the health benefits alone returned about €3 to €4 in value to society (Rojas-Rueda2011).

% It is worth noting that not all economic impacts are positive: bike-sharing has faced issues of vandalism and theft (leading to costs for replacement), and some poorly planned deployments (notably the oversupply of dockless bikes in China around 2017) resulted in financial losses and waste. The “bubble” of privately funded dockless startups in China burst, with several companies collapsing – this can be seen as a negative economic impact (capital losses, job losses) when the market failed to sustain so many operators. Nonetheless, those instances have provided lessons, and many surviving systems have adjusted to more stable business models.

% In summary, the economic impacts of BSS are multifaceted. While most systems require subsidy and are not profit-generating in themselves, their societal benefits – especially via improved health – tend to outweigh their costs when rigorously evaluated (Ricci2015; Otero2018). Users enjoy personal savings in money or time, cities gain value from healthier and potentially more mobile residents, and there are ancillary boosts to local economies and job creation. Bike-sharing should thus be viewed as a public service that delivers economic value indirectly, rather than as a stand-alone commercial venture.

% Copiat literal de Teixeira2021
% 3.4. Economy
% Modal shifts to BSS could enable significant economic gains through reducing traveltimes. Jensen, Rouquier, Ovtracht, and Robardet (2010), analysing data from Velo’v(Lyon’s BSS), found 68.2% of the system’s trips to be shorter than car trips leading toan average distance reduction of 13%. Furthermore, users of Velo’vwere reachingaverage speeds of 14.5 km/h in morning rush, close to the average car speeds in Europeanurban cores (Jensen et al., 2010). Similarly, Woodcock et al. (2014) estimated an average20% reduction in travel times as a result of shifting to the London’s BSS. Likewise, Faghih-Imani, Anowar, Miller, and Eluru (2017), comparing O/D pairs and travel times betweenBSS and taxis in New York City, discovered BSS to be either faster or competitive withtaxi for more than half of the O/D pairs within distances less than 3 km.One of the most comprehensive economic appraisals of bikesharing was conducted byBullock, Brereton, and Bailey (2017). Using survey data from Dublin’s BSS (Dublinbikes),they estimated the economic impacts from modal shift translated into environmentaland health gains, time savings as well as other wider economic benefits (WEBs) such asagglomeration gains. Even though the environmental and health gains were low dueto most of Dublinbikes’modal shift arising from walking (88.5% in this study), considerableeconomic gains deriving from time savings and WEBs were quantified. In fact, shiftingfrom walking to BSS provided significant trip time savings leading to connectivity andaccessibility improvements between home/PT stations and workplaces, which in turn pro-vided substantial agglomeration benefits. Bullock et al. (2017) then conducted a cost–benefit analysis of Dublinbikes revealing the benefits to substantially outweigh thecapital and operational costs (6.4:1 without WEBs to 12.3:1 considering WEBs), makingit competitive with major PT investments.Bikesharing could also mitigate traffic congestion. M. Wang and Zhou (2017) measuredits effect on traffic congestions in 96 US urban areas by calculating differences over time(before and after the implementation of bikesharing) and between urban areas with andTRANSPORT REVIEWS 335
% without BSS, using yearly congestion statistics. M. Wang and Zhou (2017) discovered thatin cities without BSS a 1% growth in their population was associated with a 0.6863%increase in congestion, whereas in cities with BSS a 1% increase in the population wasassociated with 0.0264% less increase in congestion, with a more pronounced effect inlarger cities. Hamilton and Wichman (2018) analysed the effect of Washington D.C.’sBSS on congestion by comparing the observed speed to a standard speed in each roadsegment at the scale of U.S. census block groups. The study reported that the presenceof BSS stations in a neighbourhood could lead to reductions in traffic congestion of upto 4% (Hamilton & Wichman, 2018). These congestion reductions may also suggest amodal shift from car to bikesharing, though none of the studies provide direct evidenceto support this supposition. Additionally, such aggregated studies are quite prone topotential confounding variables like the introduction of cycling infrastructure whichoften accompanies the implementation of BSS (Pucher, Dill, & Handy, 2010).Bikesharing may boost local businesses as explored by Buehler and Hamre (2016)inWashington D.C. Through surveying BSS users and local businesses within 0.1 miles ofBSS stations, the study found that 70% of businesses felt bikesharing to have had a posi-tive impact in the area, with 59% supporting an expansion of the system and 20% report-ing a positive impact on sales. Reciprocally, 23% of BSS users reported spending more as aresult of using bikesharing (Buehler & Hamre, 2016). Likewise, assessing the BSS of Min-neapolis/Saint Paul, X. Wang, Lindsey, Schoner, and Harrison (2015) identified that thepresence of retail businesses was positively associated with BSS trips (for each food-related business, a 1.7% increase in BSS station’s trips was observed). However, moreaccurate metrics, such as comparing retail sales within BSS service areas against overallcity trends, are needed to fully assess these possible effects.BSS also seem to influence housing prices. Controlling for spatial and temporal factors,El-Geneidy, van Lierop, and Wasfi(2016) were able to reveal a 2.7% increase in housingprices in neighbourhoods with BIXI stations. This is supported by Pelechrinis, Zacharias,Kokkodis, and Lappas (2017), which analysed the impact of Pittsburgh’s BSS onhousing and rental values by conducting quasi-experimental techniques and consideringboth intra-city and inter-city effects. Pelechrinis et al. (2017) revealed an increase in bothsales and rental prices of houses located near BSS stations, with significant increases bothat the neighbourhood and city levels. While these increases in property values can help toregenerate urban areas, they can also exacerbate inequalities if not carefully considered,potentially leading to gentrification processes, a trend already observed with othercycling policies, namely cycling lanes (Stehlin, 2015).These appraisals are of special importance as the monetary quantification of the BSS’benefits can not only help to justify public funding (Bullock et al., 2017; Chow & Sayarshad,2014; Yu et al., 2018), but also to gain support from businesses and residents, particularlywhen the implementation of BSS involves removing car space (e.g. replacing car parkingspots with BSS stations).

\subsection{Social impacts}

    % Bauman2017 té una part de socio

    % Ricci2015 diu algo de que el awarness de ciclistes puja gràcies a BSS 

% A consistent finding across studies is that \acp{bss} appeal to a specific user profile. Evidence from North America and Europe in the 2010s showed that users were mainly white, employed males~\cite{Ricci2015}, who were younger, more affluent, and more likely to cycle than the general population. This pattern has been echoed in subsequent research. Fishman et al. ~\cite{Fishman2016} noted that bike-share users are, on average, more highly educated and of higher income, and that men cycle more than women. However, it is also noted that although there is gender disparity in bike-share, it tends to be smaller than in private cycling. In low-cycling countries such as the UK, USA, and Australia, women’s participation remains lower, with studies reporting less than 20\% of trips by women in London~\cite{Goodman2014}, 22\% in Dublin~\cite{Murphy2015}, and between 23\% and 40\% in Australian schemes~\cite{Fishman2014}. By contrast, in the Netherlands, women cycle more than men~\cite{Harms2013}. 


% Recent studies confirm that these socio-demographic skews persist. A nationwide analysis of 73 U.S. cities found that bike-share users are still predominantly young, well-educated, employed males in middle- and upper-income households without children~\cite{Jin2024}. Similarly, equity assessments highlight persistent spatial imbalances, with stations disproportionately located in affluent and central neighborhoods, leaving many disadvantaged communities underserved~\cite{DuranRodas2020, Jin2024}. Research in Spain shows that gender also shapes willingness to pay and usage of e-bike sharing, with women placing greater emphasis on safety and affordability~\cite{Rodriguez2025}, while in smaller systems female users, though fewer, tend to make longer trips than men and older adults (60–69) are more active participants than previously assumed~\cite{CortezOrdonez2023}. These findings suggest that while bike-share may somewhat narrow the gender gap compared to private cycling, women, older adults, low-income groups, and minority communities remain underrepresented, with barriers such as safety concerns, station distribution, and affordability continuing to limit inclusivity.



% Equity issues extend to the spatial distribution of bike-share infrastructure. Research has found that docking stations or service areas often are concentrated in central, affluent neighborhoods, with relatively poorer coverage in low-income or peripheral areas (Duran-Rodas2020). For instance, a spatial fairness assessment of several cities’ BSS indicated imbalances in station placement that favor high-density, higher-income districts, underscoring “how fair (or unfair) the allocation of bike-sharing infrastructure” can be (Duran-Rodas2020). Such inequitable deployment may limit access for disadvantaged groups who could benefit from affordable mobility. These findings have prompted calls for more inclusive planning of BSS – e.g. offering subsidized memberships, expanding stations to underserved areas, and addressing barriers like credit card requirements – to ensure broader social inclusion in bike-sharing programs (Ricci2015). On a positive note, there is some evidence that bike-share programs can indirectly influence societal perceptions of cycling: a study in Dublin found that over 90\% of surveyed BSS users became more aware of cyclists when they themselves drove a car, suggesting bike-sharing may help improve drivers’ mindfulness of cyclists on the road (Murphy2015, as cited in Ricci2015). In summary, while BSS have provided a new mobility resource and convenience for many urban residents, the social impact has been uneven. Enhancing equity and accessibility remains a key challenge, requiring deliberate policy measures so that the benefits of bike sharing extend to a more diverse cross-section of society.

% Jin \& Sui (2024) found that disadvantaged communities with a higher share of females, youth, seniors, or minorities tend to have lower access to bike-sharing infrastructure, suggesting that users in those more privileged communities are more likely to benefit.

% Copiat literal de Bonilla-Alicea2019 pg 139
% DC bike-share reported that adding smart bike systems did not reduce smart dock system ridership, but by incorporating both systems, they increased the overall diversity of the BSS user population (Smith, 2018). A Virginia Tech Survey of Washington, DC smart dock and smart bike BSS users revealed that smart bike users were more racially diverse, had a higher proportion of female riders, and had a lower income (Virginia Tech, 2018).

% Cunha2022 - Equity impacts of cycling: examining the spatial-social distribution of bicycle-related benefits

%Copiat literalment de Teixeira2021
% 3.3. Social equityVirtually in any country, the most social disadvantaged groups tend to suffer from trans-port poverty, having lower car ownership rates and being served (or not served at all) byinefficient PT systems, which hinders their ability to access job opportunities, educationand healthcare (Lucas, Mattioli, Verlinghieri, & Guzman, 2016). Bikesharing can amelioratethese problems by providing an affordable and flexible transport alternative (Shaheenet al., 2010). Furthermore, by incentivizing physical activity, BSS can deliver additionalgains for the most social disadvantaged, as these groups tend to have higher obesityrates (Kretman Stewart, Johnson, & Smith, 2013; Xu, 2019).However, equity assessments from the UK and North America revealed that, in additionto gender unbalances (men are much more likely to use BSS than women in thosecountries), the most disadvantaged groups (including low-income, less educated andethnic minorities) are systematically under-represented among BSS users. In London,both F. Ogilvie and Goodman (2012) as well as Goodman and Cheshire (2014) foundfemales and residents in deprived areas to be severely underrepresented among usersof London’s BSS. The system’s expansion to more deprived areas and the introductionof casual trips have ameliorated some of these inequalities, but the doubling of pricesthat occurred after the expansion may have disproportionately impacted casual tripsamong poorer residents (Goodman & Cheshire, 2014). In North America, Ursaki andAultman-Hall (2016) found statistically significant differences in race, education leveland income in Chicago, Denver, Seattle and New York City; race and income in Boston;and income in Washington D.C. (which was the most equitable). McNeil, Broach, andDill (2018) observed BSS usage rates among lower-income and minorities to be lowerthan those of higher-income white residents in Brooklyn, Chicago, and Philadelphia.The main reasons for these inequalities are the location of BSS, which tend to be con-centrated near city centres to take advantage of densification (both in population andactivities) and many of these systems requiring the use of credit/debit cards (Goodman& Cheshire, 2014; McNeil et al., 2018; F. Ogilvie & Goodman, 2012; Ursaki & Aultman-Hall, 2016). Even when equity policies are implemented, the disadvantaged groups334 J. F. TEIXEIRA ET AL.
% tend to be unaware of their existence (hindering their effectiveness), due to being muchless likely to have friends or family using the systems and having less access to the Inter-net or smartphones (McNeil et al., 2018). These findings are further validated by Ber-natchez, Gauvin, Fuller, Dubé, and Drouin (2015), which through a two-years study,correlated changes in awareness levels with the proximity to BIXI stations and levels ofeducation among the population. Bernatchez et al. (2015) found that, although awarenesslevels were increasing over time, lower education was associated with lower awarenesslevels.Nevertheless, these groups have as much interest as their counterparts in using bike-sharing (McNeil et al., 2018), even showing higher usage rates per capita than the generalpopulation (F. Ogilvie & Goodman, 2012). Therefore, the more disadvantaged groupscould potentially increase their BSS usage rates with proper policy support, namely byexpanding the systems to deprived areas while maintaining them affordable and comple-menting it with target marketing campaigns (Goodman & Cheshire, 2014; McNeil et al.,2018; F. Ogilvie & Goodman, 2012; Ursaki & Aultman-Hall, 2016). This, in turn, wouldimprove the contribution of bikesharing in decreasing the current accessibility gapfound between car-owning and car-less groups (Martens, Golub, & Robinson, 2012)
